%% SMAUG-T smaug1 chosen-CT SCA — 한글 paper draft
%% main2_v1.tex 양식 차용 + kotex 추가
\documentclass{article}

\usepackage{neurips_2024}            % NeurIPS style
\usepackage[utf8]{inputenc}
\usepackage{kotex}                   % 한글 본문 처리 (XeLaTeX 도 호환)
\usepackage[T1]{fontenc}
\usepackage{microtype}
\usepackage{graphicx}
\usepackage{subcaption}
\usepackage{booktabs}
\usepackage{hyperref}
\usepackage{amsmath}
\usepackage{amssymb}
\usepackage{mathtools}
\usepackage{amsthm}
\usepackage[capitalize,noabbrev]{cleveref}
\usepackage{xspace}
\usepackage{multirow}
\usepackage{float}
\usepackage{algorithm}
\usepackage{algpseudocode}
\usepackage{enumitem}
\usepackage{placeins}

\renewcommand{\topfraction}{0.9}
\renewcommand{\bottomfraction}{0.8}
\renewcommand{\textfraction}{0.1}
\renewcommand{\floatpagefraction}{0.7}

%% Theorem 환경
\theoremstyle{plain}
\newtheorem{theorem}{Theorem}[section]
\newtheorem{proposition}[theorem]{Proposition}
\newtheorem{lemma}[theorem]{Lemma}
\newtheorem{corollary}[theorem]{Corollary}
\theoremstyle{definition}
\newtheorem{definition}[theorem]{Definition}
\newtheorem{assumption}[theorem]{Assumption}
\theoremstyle{remark}
\newtheorem{remark}[theorem]{Remark}

%% Notation macros
\newcommand{\Zq}{\mathbb{Z}_q}
\newcommand{\Rq}{\mathcal{R}_q}
\newcommand{\Rp}{\mathcal{R}_p}
\newcommand{\Rt}{\mathcal{R}_t}
\newcommand{\HW}{\mathrm{HW}}
\newcommand{\HS}{\mathrm{HS}}
\newcommand{\sign}{\mathrm{sign}}
\newcommand{\round}[1]{\lfloor #1 \rceil}

\title{8개 트레이스로 SMAUG-T smaug1 의 비밀키 복구:\\
       Chosen-Ciphertext 부채널 공격의 단일 단계 접근}

\author{%
  익명 저자 \\
  이중 익명 심사 진행 중
}

\begin{document}
\maketitle

\begin{abstract}
본 논문은 2025년 KpqC 우승 후보로 선정된 격자 기반 KEM \textbf{SMAUG-T smaug1}
에 대한 첫 통계적 부채널 공격을 제시한다. 우리의 공격은 ChipWhisperer-Lite +
Vcc 션트 + STM32F415 환경에서 \textbf{4개의 chosen-ciphertext $\times$ N=2 =
8개의 전력 트레이스 (약 3.2초의 측정)} 만으로 long-term secret key 를 복구하며,
\textbf{무작위로 생성된 9개의 키쌍 모두에서 100\%의 정확도 (분산 0)} 를 달성한다.
이 공격은 단일 단계 온라인 절차이다: \emph{오프라인 템플릿 학습 없음}, 보조
훈련 데이터 없음, EM 프로브 없음.

세 가지 방법론적 기여가 이 결과를 가능하게 한다.
첫째, \emph{chosen-CT 컨벤션 정정} — $R = \Zq[X]/(X^n+1)$ 환의 다항식 곱 구조에서
$c_1 = \alpha \cdot X^j$ 는 $\mu'_j$ 가 아니라 $\mu'_i$ 에서 $s_{(i-j) \mod n}$
을 누설시킨다. 이로 인해 $c_1 = \alpha$ 상수 chosen-CT 한 번에 256개의 비밀
계수가 동시에 노출되며, 4개의 chosen-CT 만으로 full sk 가 가능하다 (4400배 효율).
둘째, \emph{Direct attack PoI 학습} — 기존 random-µ profile 에서 학습된
PoI 가 chosen-CT 에 transfer 되지 않는다는 부정적 결과를 정량적으로 보이고,
대신 chosen-CT 데이터 자체에서 cross-design Welch-t 로 PoI 를 직접 학습하는
방법을 제안한다. 셋째, \emph{component-specific PoI + sparse 복구} — 각
component 에 해당하는 oracle pair 만으로 학습을 격리해 노이즈를 제거하고,
HW=70 sparse ternary 제약 MAP 으로 후처리한다.

부정적 발견들 — profile transfer 의 sign 반전, multi-term TVLA 의 noise floor
한계 (N=64, 256), 'T' isolated multiplier 의 round-trip 실패 — 도 정량화하여
그림 전체를 완성한다. 본 결과는 SMAUG-T smaug1 unmasked 참조 구현에
대한 SOTA 수준의 부채널 공격이다.
\end{abstract}

%=============================================================================
\section{서론 (Introduction)}
\label{sec:intro}
%=============================================================================

양자 컴퓨터의 위협에 대비한 PQC (post-quantum cryptography) 표준화는
이미 진행 중이다. 미국 NIST 는 2024년 ML-KEM (CRYSTALS-Kyber) 을
FIPS 203 으로 표준화하였으며~\citep{nist-fips203}, 한국 KpqC 컴페티션은
2025년 \textbf{SMAUG-T} 를 격자 기반 KEM 후보로 선정했다~\citep{smaugt-v4}.
이러한 표준화 후 단계에서 가장 시급한 보안 평가는 \emph{구현 보안} (implementation
security) — 특히 부채널 분석 (side-channel analysis, SCA) 에 대한 평가이다.
ML-KEM (Kyber) 는 2017년 이후 여러 SASCA 연구 (Primas et al.~\citep{primas17},
Pessl \& Primas~\citep{pessl19}, Hamburg et al.~\citep{hamburg21},
Xu et al.~\citep{xu21}) 와 mismatch 기반 공격 (Ravi et al.~\citep{ravi20})
의 대상이 되어 왔으며, HQC 에 대해서도 TCHES~24--25 의 일련의 OT-PCA / FD
oracle 공격 (Goy et al.~\citep{goy24}, Dong \& Guo~\citep{donggu025otpca,
donggu025mvfd}) 이 발표되었다.

이에 비해 SMAUG-T 에 대한 부채널 분석 문헌은 매우 빈약하다. 우리가 조사한 한
KCI 한국어 논문~\citep{smaugt-spa-kci} 은 \emph{단순 전력 분석 (SPA)} 만을
수행했으며, chosen-CT 통계 공격이나 full key recovery 는 보고되지 않았다.
본 논문은 이 공백을 \emph{first-of-kind} 결과로 채우는 동시에 \emph{SOTA
수준의 trace 비용} 을 달성한다.

\paragraph{기여 요약.}

\begin{enumerate}[leftmargin=*]
\item \textbf{Chosen-CT 컨벤션 정정 (\cref{sec:convention}).}
   기존 chosen-CT 문헌에서 묵시적으로 가정하던 $c_1 = \alpha \cdot X^j \to
   \mu'_j$ 가 $s_j$ 를 누설시킨다는 가정이 negacyclic 환에서 잘못됨을
   수학적으로 도출하고 14개의 (j$\times\alpha$) 그룹에서 128/128 비트
   round-trip 일치로 검증한다. 정정된 매핑은 $c_1 = \alpha$ 상수 한 번에
   256개의 비밀 계수를 동시 누설시키므로 4개의 chosen-CT 로 full sk 가 가능하다.

\item \textbf{Direct attack PoI (\cref{sec:directpoi}).}
   기존 PQC SCA pipeline 인 \emph{profile $\to$ transfer to attack} 가
   SMAUG-T chosen-CT 에서 sign 반전을 일으켜 zero-baseline 미달 (56\%) 임을
   보이고, 대신 chosen-CT 데이터 자체에서 cross-design Welch-t 로 PoI 를
   학습하는 \emph{single-stage} 방법을 제안한다.

\item \textbf{Component-specific PoI + Sparse 복구 (\cref{sec:component}).}
   $R^2$ 모듈 구조의 각 component 에 oracle pair 를 따로 둠으로써 학습을
   격리하고, sparse ternary $\HS=70$ 제약 MAP (greedy + 1-swap) 으로 후처리한다.
   Cross-component 95.6\% $\pm$ 5.6\% $\to$ component-specific 100\% $\pm$ 0\%.

\item \textbf{Minimum trace cost (\cref{sec:exp}).}
   N=128 캡처를 부분 추출 (sub-sampling) 하여 N=1$\sim$128 곡선을 측정,
   N=2 ($=$ 8 traces total, 3.2초 측정) 부터 5/5 keypair 모두 100\%
   복구를 달성함을 보인다. N=512 multi-seed 와 합쳐 9 datasets 9/9 키 모두
   100\% (분산 0).

\item \textbf{부정적 발견의 정량화 (\cref{sec:discussion}).}
   profile transfer 의 sign 반전 (정량 r=$-0.086$), multi-term TVLA 의 N=64,256
   noise regime ($|t| < 4.5$), N=512 phase transition ($|t|=15.20$), 'T'
   isolated multiplier 명령의 round-trip 0/6 실패 (q-modular storage 추정) 를
   모두 보고하여 \emph{honest engineering} 측면에서 그림을 완성한다.
\end{enumerate}

%=============================================================================
\section{관련 연구 (Related Work)}
\label{sec:related}
%=============================================================================

\paragraph{ML-KEM (Kyber) 부채널 공격.}
Primas et al.~\citep{primas17} 이 NTT 다항식 곱의 butterfly 구조를 SASCA
(soft analytical SCA + belief propagation) 의 factor graph 로 모델링한 첫
연구를 제시하였다. Pessl \& Primas~\citep{pessl19} 는 1M개의 Hamming weight
템플릿을 213개로 압축하고 short-cycle reduction 으로 BP 의 수렴성을 향상시켰다.
Hamburg et al.~\citep{hamburg21} 는 chosen-CT 와 sparse 입력 트릭을 결합해
masked Kyber 구현에서 2--4 traces 만으로 long-term secret key 를 복구하였다.
Xu et al.~\citep{xu21} 는 pqm4 ARM 구현에 대한 chosen-CT SPA 를 제안하여
$-O0$ reference 에서 8 traces, $-O3$ pqm4 에서 960 traces 의 성능을 보고하였다.
Bronchain \& Cassiers~\citep{bronchain20masked} 는 마스크 적용 Kyber 의
template attack 을 분석한다.

\paragraph{HQC 부채널 공격.}
Goy et al.~\citep{goy24} (TCHES~24) 는 HQC 의 RS decoder 와 RS encoder 의
Galois field multiplication 을 BP-based SASCA 로 분석하여 codeword masking
과 fine/coarse shuffling 을 모두 통과한다. Dong \& Guo~\citep{donggu025otpca}
의 OT-PCA (TCHES~25) 는 offline templates 를 활용해 hqc-128 에서 460개의
oracle calls 만으로 키 복구를 달성하며, MV-PC + FD oracle~\citep{donggu025mvfd}
은 hqc-128 에서 2개의 oracle query 로 키 복구를 보고한다.
Dong et al.~\citep{donggu025singletrace} (eprint~25) 의 single-trace 공격은
codeword 마스킹된 HQC 에 대한 1 trace 키 복구를 달성한다.

\paragraph{SMAUG-T 부채널 분석.}
우리가 조사한 한 (KCI ART003175038)~\citep{smaugt-spa-kci} 은 SMAUG-T 의
다항식 곱셈 연산에 대한 \emph{단순 전력 분석 (SPA)} 만을 수행했으며,
시각적 패턴 식별 외의 통계적 chosen-CT 공격, 전체 키 복구, masked
또는 shuffled 구현 공격은 보고되지 않았다.

\paragraph{본 연구의 위치.}

\begin{table}[h!]
\centering
\caption{PQC 부채널 공격 비교: trace 비용 vs 오프라인 비용.}
\label{tab:related}
\small
\begin{tabular}{lllll}
\toprule
연구 & 대상 & 트레이스 & 오프라인 & 비고 \\
\midrule
Dong et al. eprint25~\citep{donggu025singletrace} & hqc-1 (codeword masked) & 1 & 작음 & single-trace \\
Hamburg et al. 21~\citep{hamburg21} & Kyber512 (masked) & 2--4 & 작음 & chosen-CT + sparse \\
Dong \& Guo MV-FD 25~\citep{donggu025mvfd} & hqc-128 & 2 & 큼 (offline templates) & — \\
\textbf{본 논문} & \textbf{SMAUG-T smaug1} & \textbf{8} & \textbf{0} & chosen-CT + direct PoI \\
Xu et al. 21~\citep{xu21} & Kyber512 (-O0 ref) & 8 & 0 & chosen-CT SPA \\
Xu et al. 21~\citep{xu21} & Kyber512 (-O3 pqm4) & 960 & 0 & 최적화된 구현 \\
Dong \& Guo OT-PCA 25~\citep{donggu025otpca} & hqc-128 & 460 & 큼 & offline templates \\
Goy et al. 24~\citep{goy24} & hqc-128 (codeword masked) & $<$ 20k EM & 작음 & RS BP-SASCA \\
\bottomrule
\end{tabular}
\end{table}

본 연구의 위치는 (i) 트레이스 측면에서 SOTA 수준 (Hamburg21, eprint25, Xu21
와 동급), (ii) 오프라인 비용이 0 인 점에서 \emph{유일하게 단일 단계
온라인}, (iii) 대상 대수 구조가 격자 기반이지만 NTT 가 아닌 Toom-Cook 인
점에서 \emph{first-of-kind} 이다.

%=============================================================================
\section{사전 지식 (Preliminaries)}
\label{sec:prelim}
%=============================================================================

\subsection{격자 기반 KEM 의 일반 구조}

격자 기반 KEM 은 일반적으로 두 단계로 구성된다.
\textbf{IND-CPA PKE} 가 다항식 환 $\Rq = \Zq[X]/(X^n+1)$ 에서 module-LWE
또는 module-LWR 의 hardness 를 기반으로 IND-CPA 보안을 제공하며,
\textbf{Fujisaki-Okamoto (FO) transform}~\citep{fo99,hhk17} 이 이를 IND-CCA
KEM 으로 lifting 한다. FO 의 핵심 단계는:

\begin{enumerate}[leftmargin=*]
\item Decapsulation 시 PKE 의 \texttt{indcpa\_dec} 를 호출하여 메시지
   $\mu'$ 을 복호화한다.
\item 복호화된 $\mu'$ 으로 \emph{re-encryption} 을 수행하여
   $\mathit{ct}' = \texttt{indcpa\_enc}(pk, \mu', \mathit{seed})$ 를 만든다.
\item 입력 ciphertext 와 $\mathit{ct}'$ 을 비교 (verify) 하고, 일치 여부에 따라
   $\mathit{cmov}$ 로 shared secret 을 결정한다.
\end{enumerate}

이 구조에서 \emph{chosen-CT 공격} 은 적이 의도적으로 만든 ciphertext 를
decapsulator 에 보내고, 디코딩된 $\mu'$ 또는 mismatch flag 가 비밀키에 의존하는
정보를 누설하는지를 관찰한다.

\subsection{ML-KEM (Kyber) — NIST FIPS 203 표준}

ML-KEM 은 module-LWE 기반 KEM 이다. 핵심 파라미터 (Kyber512):
$n=256$, $k=2$ (module rank), $q=3329$ (\emph{prime}), centered binomial
분포 $\beta_\mu$ 의 secret key, 메시지 길이 $\delta = 32$ byte.

\paragraph{핵심 연산: NTT 다항식 곱.}
$q = 3329$ 가 $2n$-th root of unity 를 갖도록 선택되어 있어
\textbf{Number Theoretic Transform (NTT)} 가 사용된다.
키 생성과 암복호화의 모든 다항식 곱이 NTT 도메인에서 점별 곱 (point-wise
multiplication) 으로 계산된다. NTT 의 \emph{butterfly} 구조 (각 단계
$\ell$ 의 $2^\ell$-radix butterfly) 는 SASCA 의 factor graph 모델에 매우
적합하다 — 이것이 ML-KEM SASCA 문헌이 풍부한 이유이다.

\paragraph{지원 단계.}
$\mathrm{Kyber}\texttt{-512}/\texttt{-768}/\texttt{-1024}$ 가 각각 AES-128/192/256
에 해당하는 보안 수준을 제공한다.

\subsection{SMAUG-T smaug1 — KpqC 2025 우승}

SMAUG-T 는 module-LWR (Module Learning With Rounding) 기반의 KEM 으로,
\emph{ML-KEM 과 다른 단단함 가정} 을 사용한다. LWR 에서는 LWE 의 small noise
대신 \emph{rounding} 을 통해 정보 손실 (lossy) 을 만든다. 이로 인해 SMAUG-T 의
설계는 다음과 같이 구별된다:

\begin{itemize}[leftmargin=*]
\item $q = 2^{10} = 1024$ — \emph{2의 거듭제곱} (NTT-friendly prime 이 아님)
\item 다항식 곱이 \textbf{Toom-Cook 4-way + Karatsuba} (NTT 안 씀)
\item secret key 가 \emph{sparse ternary} ($s \in \{-1, 0, +1\}^n$) 이며 $\HS = 70$
\item Compress 후 ciphertext modulus $p = 256, p' = 32$, 메시지 modulus $t = 2$
\end{itemize}

\paragraph{smaug1 파라미터 (\cref{tab:smaug1}).}

\begin{table}[h!]
\centering
\caption{SMAUG-T smaug1 vs Kyber512 매개변수 비교.}
\label{tab:smaug1}
\small
\begin{tabular}{lll}
\toprule
파라미터 & SMAUG-T smaug1 & Kyber512 (ML-KEM 비교) \\
\midrule
타입 & MLWR & MLWE \\
$n$ (poly degree) & 256 & 256 \\
$k$ (module rank) & 2 & 2 \\
$q$ & 1024 ($=2^{10}$) & 3329 (prime) \\
$p$ (ct mod) & 256 & — (compress) \\
$p'$ (c2 mod) & 32 & — \\
$t$ (msg mod) & 2 & 2 \\
sk 분포 & sparse ternary $\{-1,0,+1\}$, $\HS=70$ & $\beta_2$ (centered binomial $\{-2..+2\}$) \\
다항식 곱 & Toom-Cook 4-way + Karatsuba & NTT \\
표준화 상태 & KpqC 2025 우승 & NIST FIPS 203 \\
\bottomrule
\end{tabular}
\end{table}

\paragraph{Decapsulation 핵심.}
SMAUG-T 의 \texttt{indcpa\_dec} 는 다음을 계산한다:
\begin{equation}
\mu'_i = \round{\frac{t}{p}\langle c_1, sk \rangle_i + \frac{t}{p'} c_{2,i}} \mod t
\label{eq:dec}
\end{equation}
여기서 $\langle c_1, sk \rangle$ 은 negacyclic 환 $R = \Zq[X]/(X^n+1)$
안의 내적이다. \emph{이 식이 본 논문의 모든 chosen-CT 분석의 출발점} 이다.

\subsection{부채널 분석 기초}

\paragraph{Welch t-test 기반 TVLA.}
Test Vector Leakage Assessment 는 두 클래스 (예: 고정 µ vs 임의 µ) 의
trace 평균 차이를 normalized 한 t-statistic 으로 측정한다.
sample 수 $T$ 의 multi-comparison 보정을 고려한 noise floor 는
$\sqrt{2 \ln T} \approx 4.5$ 이며, $|t| > 4.5$ 인 sample 을 \emph{leaky point}
로 분류한다.

\paragraph{Profile-based vs profile-free 분류기.}
\emph{Profile-based} (template) 공격은 알려진 sk 위에서 학습 단계 (offline)
와 알려지지 않은 sk 위 공격 단계 (online) 를 분리한다. \emph{Profile-free}
공격은 attack 데이터 자체에서 정보를 추출한다. 본 논문은 profile-free 의
극단적 사례이다.

\paragraph{Chosen-CT 공격 framework.}
적이 만든 ciphertext 를 decapsulator 에 보내고 다음 중 하나를 관찰:
(i) µ' 의 비트 (FD oracle), (ii) mismatch flag (PCO oracle), (iii) trace
(SCA oracle). 일반적으로 sk 의 한 비트당 1 query 이상을 요구하지만, 본 논문의
컨벤션 정정 결과는 \emph{1 query 로 256 비트 동시 노출} 을 가능케 한다.

\subsection{측정 setup}

본 연구의 측정 환경:

\begin{itemize}[leftmargin=*]
\item ChipWhisperer-Lite (CW1173, NewAE Technology), Vcc 션트 + 25\,dB LNA
\item CW308T-STM32F4 도터보드 (STM32F415RGTx), HCLK $= 7.38$\,MHz
\item ADC samples per trace: 24\,400 at clkgen\_x4 ($= 29.5$\,MS/s)
\item 펌웨어: SimpleSerial v1.1 기반, SMAUG-T smaug1 ROM 33\,683 byte
\item Throughput: $\mathrm{D} \approx 0.77$ tr/s (full kem\_dec),
   $\mathrm{Z} \approx 2.65$ tr/s (indcpa\_dec only).
\end{itemize}

펌웨어는 다음 SimpleSerial 명령을 노출한다: \texttt{F} (지속 keygen),
\texttt{I} (chosen-CT chunk inject), \texttt{L} (무결성 fingerprint),
\texttt{M} (PKE-labeled enc), \texttt{D} (full kem\_dec),
\texttt{Z} (indcpa\_dec only $+$ 32B µ' 응답), \texttt{X} (sk PKE dump,
ground truth verification only).

%=============================================================================
\section{Chosen-CT 컨벤션 정정 (Convention Correction)}
\label{sec:convention}
%=============================================================================

\subsection{원래 가정 (잘못됨)}

기존 SMAUG-T chosen-CT 공격 설계 시도에서 다음과 같은 매핑을 묵시적으로 가정해
왔다:
\begin{quote}
\textbf{(가정 — 잘못됨)} $c_1 = \alpha \cdot X^j$ chosen-CT 는 $\mu'_j$ 에서
$s_j$ 를 누설시킨다.
\end{quote}
이 가정 하에서는 256개의 위치 $j$ 마다 별개의 chosen-CT 가 필요하며 (oracle
pair $\alpha \in \{64, 192\}$ 와 component $\in \{0,1\}$ 적용 시) 약 $32$k 이상의
trace 가 요구된다.

\subsection{수학적 정정}

$R = \Zq[X]/(X^n+1)$ 안의 negacyclic 환 곱:
\begin{equation}
\langle c_1, sk \rangle(X) = \alpha \cdot X^j \cdot sk(X) =
   \alpha \sum_{k=0}^{n-1} s_k \cdot X^{j+k}.
\label{eq:rawmul}
\end{equation}
$X^n + 1$ 으로 reduction 하면 (anticyclic wrap, sign flip):

\begin{proposition}[chosen-CT 매핑 정정]
\label{prop:convention}
$c_1 = \alpha \cdot X^j$, $c_2 = 0$ 인 chosen-CT 의 decapsulation 결과 $\mu'$
의 $i$ 번째 비트는 다음을 만족한다:
\begin{equation}
\mu'_i = \round{\frac{t}{p} \alpha \cdot s_{(i-j) \mod n} \cdot \sigma_l(i)} \mod t
\label{eq:correct-mapping}
\end{equation}
여기서 $\sigma_l(i) = +1$ if $i \geq j$ else $-1$ 은 negacyclic wrap 부호이다.
\end{proposition}

즉, $\mu'_i$ 는 $s_j$ 가 아니라 $s_{(i-j)}$ 를 누설시킨다. 우리가 $\mu'_j$
한 비트만을 본다면 (즉 $i = j$ 이면 $i - j = 0$) \emph{모든 $j$ 에서 동일하게
$s_0$ 만 측정} 된다 — 이것이 14개의 (j$\times\alpha$) 그룹 pilot 캡처에서
"왜 모든 $j$ 가 같은 답을 주는가?" 라는 의문을 낳은 결정적 단서였다.

\subsection{새로운 chosen-CT 디자인}

식~\eqref{eq:correct-mapping} 을 $j = 0$ 에 적용:
\begin{equation}
\mu'_i = \round{\frac{t}{p} \alpha \cdot s_i} \mod t \quad \forall i \in [0, n).
\label{eq:const-c1}
\end{equation}

즉 \textbf{$c_1 = \alpha$ (상수 다항식, $j=0$) 인 \emph{단 하나의} chosen-CT 가
모든 256개의 비밀 계수 $s_0, \ldots, s_{255}$ 를 동시에 $\mu'$ 으로 누설시킨다}.

oracle pair $(\alpha_{\mathrm{pos}}=64, \alpha_{\mathrm{neg}}=192)$ 와 module rank
$k=2$ 의 두 component 를 결합하면 \textbf{4개의 chosen-CT} 로 full sk 가 가능하다:

\begin{table}[h!]
\centering
\caption{4개의 chosen-CT 디자인.}
\label{tab:4ct}
\small
\begin{tabular}{lll}
\toprule
chosen-CT 라벨 & $c_1$ & 측정하는 정보 \\
\midrule
$D_0^+$ & $c_1[0] = 64,\; c_1[1]=0$ & component 0 의 $+1$ detector \\
$D_0^-$ & $c_1[0] = 192,\; c_1[1]=0$ & component 0 의 $-1$ detector \\
$D_1^+$ & $c_1[0] = 0,\; c_1[1]=64$ & component 1 의 $+1$ detector \\
$D_1^-$ & $c_1[0] = 0,\; c_1[1]=192$ & component 1 의 $-1$ detector \\
\bottomrule
\end{tabular}
\end{table}

\subsection{Multi-term 과 $R^2$ cross-component 일반화}

일반화된 multi-term chosen-CT 정의:
$$
c_1[m] = \sum_l \alpha_{m,l} \cdot X^l \quad (m \in \{0, 1\}, l \in [0, n)).
$$
\begin{proposition}[일반화 매핑]
\label{prop:multi-term}
$c_1$ 의 일반 형태에서 decapsulation 결과의 $i$ 번째 계수는:
\begin{equation}
\mu'_i = \round{ \frac{t}{p} \sum_{m, l} \alpha_{m, l} \cdot s_m^{(i-l) \mod n}
   \cdot \sigma_l(i) + \frac{t}{p'} c_{2,i} } \mod t
\label{eq:multi-term}
\end{equation}
\end{proposition}
이 일반 매핑은 $c_1[0] = 64 \cdot X^l, c_1[1] = 64 \cdot X^k$ 같은
\emph{$R^2$ cross-component} 디자인 (eprint~25 OT-FDA-CK 영감)~\citep{donggu025singletrace}
도 포함한다.

\subsection{보드 round-trip 검증}

본 논문의 \texttt{host/smaug/chosen.py::predict\_mu\_prime} 가 식~\eqref{eq:multi-term}
을 구현하며, 7개의 다양한 design (\texttt{H\_const\_a64, H\_const\_a192,
H\_2term, H\_3term, H\_combined\_l10\_k50, ...}) 의 보드 'Z' 응답과 비교 시
\textbf{모든 256/256 비트 일치 (round-trip 100\%)}.
또한 14개의 (j $\times \alpha$) pilot 그룹의 Z 응답을 'X' sk dump 와
식~\eqref{eq:correct-mapping} 으로 재해석할 때 \textbf{128/128 비트 일치}.
이 검증은 \texttt{tests/test\_chosen.py} (15 tests) + \texttt{test\_partition.py}
(14 tests) 에 자동화되어 있다.

%=============================================================================
\section{Direct Attack PoI 학습}
\label{sec:directpoi}
%=============================================================================

\subsection{Profile-PoI 의 한계 (negative result)}

표준 PQC SCA pipeline 은 두 단계이다: (i) \emph{profile} 단계에서 알려진 sk
위에서 random µ 의 trace 를 모아 per-bit Welch t-test 로 PoI[$i$] 와 sign[$i$]
를 학습, (ii) \emph{attack} 단계에서 알려지지 않은 sk 위에서 chosen-CT 의
trace 에 학습한 PoI 를 적용. 우리는 이 pipeline 을 1000개의 random µ profile
trace 로 시도하였다.

\paragraph{결과 (negative).} bit accuracy 56\% — \emph{zero-baseline (모두 0
예측) 73\%에 미달}. 자세한 진단:
\begin{itemize}[leftmargin=*]
\item $d_i$ (signed score) 와 ground-truth $sk_i$ 의 correlation $r = -0.086$
   (\emph{부호 반전!}).
\item Sign flip 후 sign accuracy 35\% $\to$ 65\% 향상되지만 bit accuracy 는
   여전히 chance level.
\end{itemize}

원인 분석: profile (random µ, binomial(0.5)) 와 chosen-CT (deterministic µ
per design) 의 trace 가 서로 다른 \emph{baseline power offset} 을 갖는다.
Profile 에서 학습된 sign 은 \emph{cross-domain transfer} 에서 안정적이지
않다.

\subsection{Direct attack PoI: 핵심 method}

핵심 통찰: 각 chosen-CT design 의 µ' 은 \emph{deterministic per-design} 이며
(같은 sk + 같은 input → 같은 µ'), 4개의 design 이 모이면 각 비트 $i$ 에 대해
서로 다른 0/1 split 을 만든다. 이 \emph{cross-design split} 위에서 직접
Welch-t 를 계산할 수 있다.

\begin{algorithm}[h!]
\caption{Direct attack PoI 학습 (component-specific)}
\label{alg:directpoi}
\begin{algorithmic}[1]
\Require attack traces $T \in \mathbb{R}^{N \times S}$, µ' bits per trace $B \in \{0,1\}^{N \times 256}$,
   trace label $\mathit{comp} \in \{0, 1\}^N$, target component $c$, min class size $m$
\Ensure PoI $p \in \{-1\} \cup [0, S)^{256}$, sign $\sigma \in \{-1, +1\}^{256}$
\State $T_c \gets T[\mathit{comp} = c]$,\; $B_c \gets B[\mathit{comp} = c]$
\For{$i = 0$ to $255$}
\State $n_0 \gets |\{j: B_c[j, i] = 0\}|$,\; $n_1 \gets |\{j: B_c[j, i] = 1\}|$
\If{$n_0 < m$ or $n_1 < m$}
   \State $p[i] \gets -1$ \Comment{invalid bit (insufficient cross-design split)}
\Else
   \State $\mu_0 \gets \mathrm{mean}(T_c[B_c[:, i] = 0]),\; \mu_1 \gets \mathrm{mean}(T_c[B_c[:, i] = 1])$
   \State $\mathrm{Var}_0, \mathrm{Var}_1 \gets$ pooled variances
   \State $t[\cdot] \gets (\mu_1 - \mu_0) / \sqrt{\mathrm{Var}_0 / n_0 + \mathrm{Var}_1 / n_1}$
   \State $p[i] \gets \arg\max_s |t[s]|,\; \sigma[i] \gets \mathrm{sgn}(t[p[i]])$
\EndIf
\EndFor
\Return $(p, \sigma)$
\end{algorithmic}
\end{algorithm}

\subsection{왜 valid bits = sk nonzero positions 인가?}

oracle pair $(\alpha=64, \alpha=192)$ 에서 식~\eqref{eq:const-c1} 의 round
함수에 의해:
\begin{itemize}[leftmargin=*]
\item $s_i = 0 \Rightarrow$ 두 design 모두 $\mu'_i = 0$ (cross-design split 비어
   있음, invalid bit).
\item $s_i = +1 \Rightarrow \alpha=64$ 가 $\mu'_i = 1$, $\alpha=192$ 가 $\mu'_i = 0$.
\item $s_i = -1 \Rightarrow$ 반대.
\end{itemize}
즉 \emph{Direct attack PoI 의 valid bit 집합은 정확히 $sk$ 의 support
(nonzero positions)} 이며, 이는 $\HS = 70$ 이다. 따라서 sparse ternary
$\HS = 70$ 제약은 \emph{디자인 자체에 implicit} 으로 들어 있다.

%=============================================================================
\section{Component-specific PoI 와 Sparse 복구}
\label{sec:component}
%=============================================================================

\subsection{Cross-component 노이즈 제거}

4개의 chosen-CT 모두에서 PoI 를 학습하는 \emph{cross-component} 방식은
component 1 design 의 trace 가 component 0 의 PoI 학습에 노이즈를 추가한다
(반대도 마찬가지). 실제 측정에서 N=128 cross-component PoI 의 평균 정확도는
95.6\% $\pm$ 5.6\% 이며 worst-case seed 가 87.1\% 까지 떨어진다.

\paragraph{개선 — Component isolation.} $s_c$ 를 공격할 때, label$\_$comp == $c$
인 두 oracle pair 만 사용해 PoI 를 학습한다. 이 isolation 으로 N=128 평균
정확도가 \textbf{100\% $\pm$ 0\% (5/5 seeds)} 로 향상된다.

\subsection{Sparse Ternary 복구 (HW=HS MAP)}

$d_i = \sigma_i \cdot (\bar{T}_{\alpha=64}[\mathrm{PoI}_i] -
\bar{T}_{\alpha=192}[\mathrm{PoI}_i])$ 의 signed score 를 다음과 같이 ternary
posterior 로 변환:
\begin{equation}
\pi_i(s) \propto \exp\left( -\frac{1}{2 \sigma^2} (d_i - \mu \cdot s)^2 \right)
\quad s \in \{-1, 0, +1\}
\end{equation}
여기서 $\mu$ 는 $|d|$ 의 추정 magnitude.

\begin{algorithm}[h!]
\caption{Greedy sparse ternary recovery (HW=HS MAP)}
\label{alg:sparse-recover}
\begin{algorithmic}[1]
\Require posterior $\pi \in \Delta^{n \times 3}$, sparse constraint $\HS$
\Ensure $\hat{s} \in \{-1, 0, +1\}^n$ with $|\{i: \hat{s}_i \neq 0\}| = \HS$
\State $r_i \gets \pi_i(-1) + \pi_i(+1)$ for all $i$ \Comment{support score}
\State $\mathit{Top} \gets \arg\mathrm{topk}(r, \HS)$ \Comment{stable sort, top-HS positions}
\State $\hat{s} \gets \mathbf{0}$
\For{$i \in \mathit{Top}$}
   \State $\hat{s}_i \gets +1$ if $\pi_i(+1) > \pi_i(-1)$ else $-1$
\EndFor
\Return $\hat{s}$
\end{algorithmic}
\end{algorithm}

\paragraph{1-swap neighborhood.} 정확도가 부족한 경우 top-$\HS$ 결과의
1-swap 변종 (in-support 하나를 빼고 out-of-support 하나를 대체) 을 enumerate
하고 log-prob 정렬한다. 이 방법은 \texttt{enumerate\_top\_candidates} 로
구현되어 있다.

\paragraph{$\sigma$ 정확도.} 4 chosen-CT design 의 cross-design split 이 sk
support 와 정확히 일치하는 점 (\cref{sec:directpoi}) 때문에, sparse\_recover
의 \emph{선택된 support 안에서 sign 정확도는 9개 keypair 모두에서 100\%}
이다.

%=============================================================================
\section{실험 평가 (Experiments)}
\label{sec:exp}
%=============================================================================

\subsection{측정 setup 요약}

ChipWhisperer-Lite + Vcc 션트 + STM32F415, ADC 24\,400 samples,
clkgen\_x4 (29.5\,MS/s), gain 25\,dB. 펌웨어는 SimpleSerial v1.1, 'Z'
명령으로 trace + 32 byte $\mu'$ 응답 수신. 자세한 명령 셋과 환경은
\cref{sec:prelim}, repository \texttt{README.md} 참조.

\subsection{N curve — minimum trace cost}

\begin{table}[h!]
\centering
\caption{N curve (5 seeds × N=128 캡처를 부분 추출). 4 chosen-CT 의 각각
처음 $N$ trace 만 사용. Phase transition 이 $N=1 \to N=2$ 에서 발생.}
\label{tab:ncurve}
\small
\begin{tabular}{rrrrr}
\toprule
$N$ (per CT) & total traces & 캡처 시간 & full sk acc (5 seeds) & std \\
\midrule
1 & 4 & 1.6 s & 57.3\% & 0.015 \\
\textbf{2} & \textbf{8} & \textbf{3.2 s} & \textbf{100\%} & \textbf{0\%} \\
4 & 16 & 6.4 s & 100\% & 0\% \\
16 & 64 & 25.6 s & 100\% & 0\% \\
128 & 512 & 204.8 s & 100\% & 0\% \\
\bottomrule
\end{tabular}
\end{table}

$N=1$ 은 Welch t-statistic 의 variance 추정 (ddof=1) 에 최소 2 sample 이
필요하기 때문에 실패한다. $N \geq 2$ 부터 모든 5개의 keypair 에서 full sk
복구가 100\% 달성된다.

\subsection{9-seed grand evaluation}

5 seeds $\times$ N=128 와 3 seeds $\times$ N=512 와 1 seed $\times$ N=256 을
합쳐 \textbf{9 dataset 평가}:

\begin{table}[h!]
\centering
\caption{9 dataset multi-seed 평가 (component-specific PoI + sparse\_recover).}
\label{tab:9seed}
\small
\begin{tabular}{lrrrrr}
\toprule
metric & mean & std & min & max & median \\
\midrule
\textbf{full\_sk\_acc} & \textbf{1.000} & \textbf{0.000} & \textbf{1.000} & \textbf{1.000} & \textbf{1.000} \\
n\_valid\_bits & 120.3 & 2.8 & 116 & 125 & 122 \\
$\max|t|$ & 10.56 & 3.04 & 6.7 & 15.3 & 9.2 \\
$s_0$ corr & +1.000 & 0.000 & +1.000 & +1.000 & +1.000 \\
$s_1$ corr & +1.000 & 0.000 & +1.000 & +1.000 & +1.000 \\
\bottomrule
\end{tabular}
\end{table}

\subsection{SNR phase transition (negative finding)}

\paragraph{Cross-component PoI 의 한계.} component 격리 없이 4 design 모두를
사용한 PoI 학습:

\begin{table}[h!]
\centering
\caption{Cross-component PoI 의 SNR phase transition.}
\label{tab:snr-phase}
\small
\begin{tabular}{lrrr}
\toprule
비교 & $N=64\;\max|t|$ & $N=256\;\max|t|$ & $N=512\;\max|t|$ \\
\midrule
const\_a64 vs const\_a192 & — & — & 13.88 \\
const\_a64 vs 3term & 4.03 & 3.64 & 15.20 \\
const\_a64 vs combined & — & — & 13.80 \\
\bottomrule
\end{tabular}
\end{table}

$N=64 \to N=256$ (4배 증가) 에서 $|t|$ 가 4.03 $\to$ 3.64 로 \emph{역행}
(noise regime, random walk). $N=256 \to N=512$ (2배) 에서 $|t|$ 가
3.64 $\to$ 15.20 (4.2배 증가, $\sqrt{2} \approx 1.41$ 보다 큼) — \emph{noise
$\to$ signal regime phase transition}.

\subsection{비교 method 와 차이}

\begin{table}[h!]
\centering
\caption{Cross-component PoI vs Component-specific PoI (5 seeds, N=128).}
\label{tab:methods}
\small
\begin{tabular}{lrrrr}
\toprule
metric & Cross-component & \textbf{Component-specific (paper main)} \\
\midrule
full sk mean & 95.6\% & \textbf{100\%} \\
full sk std & 5.6\% & \textbf{0\%} \\
full sk min & 87.1\% & \textbf{100\%} \\
$s_0$ mean & 98.8\% & \textbf{100\%} \\
$s_1$ mean & 92.5\% & \textbf{100\%} \\
sign acc & 100\% & 100\% \\
\bottomrule
\end{tabular}
\end{table}

%=============================================================================
\section{토의 (Discussion)}
\label{sec:discussion}
%=============================================================================

\subsection{무엇이 작동했고 무엇이 안 했나}

\paragraph{성공 요인.}
(i) chosen-CT 컨벤션 정정으로 query 효율 4400배 향상. (ii) Direct attack PoI
가 cross-domain transfer 문제 우회. (iii) Component-specific learning 이
cross-component 노이즈 제거. (iv) sparse\_recover 의 $\HS$ 제약이 design
자체에 implicit.

\paragraph{실패 시도 (negative results, 정량).}
\begin{enumerate}[leftmargin=*]
\item \textbf{Random-µ profile transfer} (\cref{sec:directpoi}). Sign 반전,
   bit accuracy 56\% (zero-baseline 73\% 미달). $r=-0.086$, sign flip 후
   $r=+0.086$ 으로 호전되나 여전히 chance.
\item \textbf{Per-bit FD attack at noise floor.} Profile $N=200 \to 1000$ (5배)
   에서 max$|t|$ 가 4.97 $\to$ 5.19 — $\sqrt{5} \approx 2.24$ 의 기대 증가량
   (예측 11.1) 보다 훨씬 작음, noise peak 만 측정.
\item \textbf{'T' isolated multiplier 명령 round-trip.} $\texttt{poly\_mul\_acc}$ 를
   firmware 의 'T' 명령으로 isolated 호출 + sub-trigger. host 의 numpy 시뮬레이션과
   board 의 응답이 \emph{0/6 round-trip 매치}. 패턴 분석:
   board=$0\mathrm{x}7\mathrm{FC}0$ vs expected $0\mathrm{xFFC}0$ — high bit
   차이. SMAUG-T archive 가 q-modular unsigned (또는 NTT-friendly variant) 의
   internal storage 사용 추정. 이 issue 우회를 위해 'Z' 명령 (전체
   indcpa\_dec) trace windowing 으로 pivot.
\end{enumerate}

\subsection{Threat model 정직 평가}

본 공격은 다음 능력을 가진 적을 가정한다:
\begin{itemize}[leftmargin=*]
\item Decapsulator 에 임의의 ciphertext 를 보낼 수 있음 (chosen-CT capability)
\item Decapsulation 동안 Vcc 션트 또는 EM 위에서 power trace 측정 가능
\item Trace 측정 transparency (트리거 동기화) — 일반적인 SCA 가정
\end{itemize}
적은 \emph{오프라인 템플릿} 을 학습할 필요가 \emph{없으며}, target sk 의 다른
복사본이나 alternative 정보 통로도 필요 없다. 이는 \emph{single-stage online}
공격이며 PCO/MV-PC 등의 offline-heavy 공격과 차별점이다.

%=============================================================================
\section{대응 방안 분석 (Countermeasures)}
\label{sec:countermeasures}
%=============================================================================

본 공격에 대한 표준 PQC SCA 대응 방안의 정량적 효과를 분석한다.
실측 구현은 향후 연구로 남기되, 기존 문헌 (HQC TCHES~24 의 shuffling
breakdown, masked Kyber 평가 결과) 에 기반한 분석은 다음과 같다.

\subsection{Masking}

\paragraph{1차 masking (codeword 또는 arithmetic).}
$\langle c_1, sk \rangle_i$ 를 두 share 로 분할 ($\mathrm{share}_a + \mathrm{share}_b
\equiv \langle c_1, sk \rangle_i \mod q$). $\mu'$ 응답 자체는 unmask 후 출력되어
변하지 않으나, trace-level signal 은 두 share 의 incoherent 합으로 약 $\sqrt{2}$
배 감소한다.
\textbf{예상 trace 비용}: $N=2 \to N \approx 4-8$, 즉 16--32 traces, 약 10--15초.

\paragraph{2차 masking.}
3 share 분할 후 univariate t-test 의 SNR 이 약 3--4배 감소, bivariate Welch-t
의 multi-comparison floor 가 $\sqrt{2 \ln T^2}$ 으로 증가. 추가로 sample
pair search 의 계산 비용이 $T^2$.
\textbf{예상 trace 비용}: $N \approx 16-32$.

\paragraph{3차 이상 masking.}
First-order Welch-t method 가 fail. BP-SASCA / template attack 영역으로 진입.

\subsection{Shuffling}

\paragraph{Sample-level (random delay).}
PoI 가 fixed sample 이 아닌 window 안의 임의 sample 로 spread.
window 폭 $w$ 에 대해 SNR 이 $\sqrt{w}$ 배 감소. $w=8$ 이면 $N \approx 16$
(64 traces), $w=32$ 이면 $N \approx 64$ (256 traces).

\paragraph{Coefficient-level shuffling (Ravi-Roy 2020 RPBC).}
Compress$_t$ 의 비트 처리 순서를 random permute. PoI 가 fixed sample
대응 안 됨. 강한 방어.

\paragraph{Full shuffling (Goy-Maillard 24 영감)~\citep{goy24}.}
Sample + coefficient-level 결합. \textbf{본 method 직접 fail}. Resync
+ BP-SASCA 가 필요.

\subsection{종합 권고}

\begin{table}[h!]
\centering
\caption{대응 방안 별 본 공격의 비용 추정.}
\label{tab:counter}
\small
\begin{tabular}{lrl}
\toprule
방어 & 예상 trace & 추가 method change \\
\midrule
없음 (현재) & \textbf{8} & — \\
1차 masking & 16--32 & 없음 \\
2차 masking & 64+ & bivariate Welch-t \\
Sample shuffle ($w=8$) & 64 & 없음 \\
Coefficient shuffle & 큼 & PoI rebalancing \\
\textbf{Full shuffling} & \textbf{fail} & resync + BP-SASCA 필요 \\
\bottomrule
\end{tabular}
\end{table}

권장 defense: \textbf{1차 masking + sample-shuffle ($w \geq 16$) +
coefficient-shuffle 결합}. 각각 individually 우회 가능하나 결합 시 본
method 가 BP-SASCA 영역으로 강제됨 (SMAUG-T 에는 prior work 없음).

%=============================================================================
\section{결론 (Conclusion)}
\label{sec:conclusion}
%=============================================================================

본 논문은 KpqC 2025 우승 KEM SMAUG-T smaug1 에 대한 first-of-kind
통계적 부채널 공격을 제시한다. \textbf{ChipWhisperer-Lite + Vcc 션트 +
STM32F415 unmasked 참조 구현 위에서 4개의 chosen-CT $\times$ N=2 = 8개의
power trace (3.2초 측정) 만으로 long-term secret key 의 100\% 복구를 9개의
무작위 keypair 모두에서 분산 0 으로 달성한다}.

세 method 기여 (chosen-CT 컨벤션 정정, Direct attack PoI, Component-specific
PoI + Sparse 복구) 와 정량화된 부정적 발견 (profile transfer fail, SNR
phase transition, 'T' isolated multiplier round-trip fail) 이 결과를 가능케
하였다. 본 method 는 module-LWR 또는 module-LWE 의 임의의 KEM 으로
일반화 가능하며, 향후 연구는 (i) masked SMAUG-T 구현 평가, (ii) smaug3 /
smaug5 (모듈 rank 3 / 5) 확장, (iii) Toom-Cook BP-SASCA 의 결합으로
single-trace ($N=1$) 가능성 확인이다.

%=============================================================================
\section*{재현성 (Reproducibility)}
\label{sec:repro}
%=============================================================================

본 논문의 모든 실험은 다음으로 재현 가능하다:

\begin{verbatim}
# 1. firmware 빌드 + 플래시
make -C firmware/simpleserial-smaug PLATFORM=CW308_STM32F4 SMAUG_LEVEL=1
python3 host/upload.py firmware/simpleserial-smaug/\
    simpleserial-smaug-CW308_STM32F4.hex

# 2. multi-seed 캡처 (5 keypairs × N=128, ~25분)
for s in 1 2 3 4 5; do
    python3 scripts/run_attack.py -n 128 \
        --out traces/attack_seed${s}.npz
done

# 3. paper main 결과 분석
python3 scripts/analyze_multi_seed.py traces/attack_seed{1..5}.npz
# 출력: full_sk_acc mean=1.000 std=0.000

# 4. minimum cost 데모 (8 traces, 3.2초)
python3 scripts/run_attack.py -n 2 --out traces/attack_quick.npz
python3 scripts/analyze_multi_seed.py traces/attack_quick.npz
\end{verbatim}

코드는 \texttt{<repository>} 에서 다운로드 가능하며, 64개의 unit test
(\texttt{tests/}) 가 chosen-CT 빌더, predict\_mu\_prime 의 round-trip,
sparse\_recover 의 합성 SNR 곡선을 자동 검증한다.

%=============================================================================
\bibliographystyle{plainnat}
\bibliography{mybib}

%=============================================================================
\appendix
%=============================================================================

\section{SMAUG-T 매개변수 상세}
\label{app:params}

\subsection{smaug1, smaug3, smaug5}

\begin{table}[h!]
\centering
\caption{SMAUG-T 모든 보안 수준의 파라미터.}
\label{tab:smaug-all}
\small
\begin{tabular}{lrrr}
\toprule
파라미터 & smaug1 & smaug3 & smaug5 \\
\midrule
$n$ & 256 & 256 & 256 \\
$k$ (module rank) & 2 & 3 & 5 \\
$\log q$ & 10 & 10 & 11 \\
$\log p$ & 8 & 8 & 8 \\
$\log p'$ & 5 & 8 & 6 \\
$\log t$ & 1 & 1 & 1 \\
$\HS$ (per poly) & 70 & 88 & 152 \\
ciphertext bytes & 672 & 1088 & 1472 \\
공개키 bytes & 672 & 992 & 1632 \\
보안 수준 (target) & AES-128 & AES-192 & AES-256 \\
\bottomrule
\end{tabular}
\end{table}

본 논문의 모든 실험은 smaug1 위에서 수행되었다. smaug3 / smaug5 의 module
rank 3 / 5 는 더 많은 chosen-CT design (module rank $\times$ 2 = 6 / 10 design)
이 필요하지만 method 자체는 동일하다.

\subsection{Toom-Cook 4-way layout (smaug1.a 분석)}

\texttt{arm-none-eabi-objdump} 으로 \texttt{lib/crypto\_kem/smaug1.a} 를 분석한
결과:

\begin{table}[h!]
\centering
\caption{핵심 함수의 instruction 수와 cycle 추정 (Cortex-M4).}
\label{tab:toom-cook}
\small
\begin{tabular}{lrrr}
\toprule
함수 & instructions & 추정 cycles & samples (clkgen\_x4) \\
\midrule
\texttt{karatsuba\_simple} (1회) & 601 & ~1500 & ~6000 \\
\texttt{toom\_cook\_4way} (자체) & 651 & ~1000 & ~4000 \\
\texttt{toom\_cook\_4way} (전체, +7 kara) & 5808 & ~11500 & ~46000 \\
\texttt{poly\_mul\_acc} & 5882 & ~12000 & ~48000 \\
\texttt{indcpa\_dec} (전체) & ~5000 & ~50000 & ~200000 \\
\bottomrule
\end{tabular}
\end{table}

\texttt{poly\_mul\_acc} 는 \texttt{toom\_cook\_4way} 를 한 번 호출하며, 그
내부에서 \texttt{karatsuba\_simple} 이 7번 호출되어 Toom-Cook-4 의 7개의
evaluation point 별 sub-poly 곱을 수행한다. 24\,400 sample ADC 윈도우는
\texttt{poly\_mul\_acc} 한 호출 (46\,000 samples) 의 절반만 포함하므로,
\texttt{Z} 명령은 \texttt{indcpa\_dec} 의 시작 영역만 캡처한다.

\section{추가 실험 디테일}
\label{app:exp}

\subsection{발견 과정 timeline}

본 논문의 main flow 는 다음 과정을 거쳐 도출되었다:

\begin{enumerate}[leftmargin=*]
\item \textbf{Step 1.} Collective µ flip leakage 측정 (E3b TVLA, $N=1000$,
   max$|t|=8.99$). \emph{paper main 에는 인용만}, 코드는
   \texttt{history/scripts/capture\_step1\_collective\_leak\_n1000.py}.
\item \textbf{Step 2.} 14개의 chosen-CT pilot 으로 컨벤션 발견 trigger
   (j-sweep 이 모두 같은 $s_0$ 만 측정함을 발견). 코드는
   \texttt{history/scripts/capture\_step2\_convention\_pilot.py}.
\item \textbf{Step 3.} Random-µ profile 캡처 후 cross-domain transfer 시도
   (negative). 코드는
   \texttt{history/scripts/capture\_step3\_random\_mu\_profile.py}.
\item \textbf{Step 4a/4b.} Profile-PoI baseline (56\%) + sparse 후처리 (55\%)
   결과 정량화 (negative).
\item \textbf{Step 5.} multi-term H 데이터의 cross-design split 첫 분석 →
   direct attack PoI 도출 (paper Section 5).
\end{enumerate}

이 과정의 각 step 은 \texttt{history/scripts/} 에 보존되어 있으며 paper 의
부정적 또는 prerequisite finding 의 evidence 역할을 한다.

\subsection{재현성 file mapping}

paper 의 각 section 은 다음 main script 로 재현된다:

\begin{table}[h!]
\centering
\caption{Paper section \ ↔ \ reproducer script mapping.}
\label{tab:repro}
\small
\begin{tabular}{ll}
\toprule
Paper Section & Reproducer \\
\midrule
Section 4 (Method overview) & \texttt{scripts/run\_attack.py} \\
Section 4.4 (Multi-term + R\textsuperscript{2} 검증) & \texttt{scripts/run\_h\_attack.py} \\
Section 5 (Direct attack PoI) & \texttt{scripts/attack\_direct\_poi.py} \\
Section 6, 7 (\textbf{paper main result}) & \textbf{\texttt{scripts/analyze\_multi\_seed.py}} \\
\bottomrule
\end{tabular}
\end{table}

\end{document}
